\begin{notechapter}
\chapter{Convergence of Events and Probabilities}\label{ch:convergence-of-events}
\chaptersubtitle{When does $\Prob(A_n)\to\Prob(A)$?}

\begin{chapterintro}
Proofs in learning theory often approximate an event by a family of simpler
events and pass to the limit. Doing so requires knowing when probabilities
follow events to the limit. This chapter answers that question and applies the
answer to the symmetrization lemma of \cref{ch:symmetrization-lemma}. You will
learn
\begin{itemize}[nosep]
  \item what it means for events to converge: $\liminf$, $\limsup$ and the
        set-theoretic limit;
  \item continuity of probability along monotone sequences, and the
        convergence theorem for arbitrary convergent sequences;
  \item convergence in symmetric difference, and Fatou's bounds;
  \item examples where events do not converge but probabilities do, and
        where probabilities do not converge;
  \item how the symmetrization proof avoids taking a limit of events.
\end{itemize}
\prerequisites{basic set operations; a probability space $(\Omega,\mathcal F,\Prob)$;
the dominated convergence theorem (for one proof).}
\end{chapterintro}

\section{Limits of events}\label{convergence-of-events:sec:limits}

Let $(A_n)_{n\geq1}$ be a sequence of events. Statements about families indexed by
a real parameter $\epsilon\to0$ reduce to sequences by taking $\epsilon_n\to0$.

\begin{definition}[Limit inferior and superior]
\[
  \liminf_nA_n = \bigcup_{N\geq1}\bigcap_{n\geq N}A_n,\qquad
  \limsup_nA_n = \bigcap_{N\geq1}\bigcup_{n\geq N}A_n .
\]
The sequence \emph{converges} to $A$ (in the set-theoretic sense) if
$\liminf_nA_n = \limsup_nA_n = A$.
\end{definition}

The definitions are easier to read outcome by outcome:
\begin{itemize}
  \item $\omega\in\liminf_nA_n$ iff $\omega$ belongs to $A_n$ for \emph{all $n$
        large enough} (``eventually always'');
  \item $\omega\in\limsup_nA_n$ iff $\omega$ belongs to $A_n$ for \emph{infinitely
        many} $n$ (``infinitely often'').
\end{itemize}
Always $\liminf A_n\subseteq\limsup A_n$. Convergence means that every outcome
is eventually always in, or eventually always out: equivalently, the indicators
converge pointwise, $\bbone_{A_n}(\omega)\to\bbone_A(\omega)$ for every
$\omega$.

\section{When probabilities converge}\label{convergence-of-events:sec:probabilities}

\begin{theorem}[Continuity of probability]\label{convergence-of-events:thm:monotone}
\begin{enumerate}[nosep]
  \item If $A_1\subseteq A_2\subseteq\cdots$, then $\Prob(A_n)\to\Prob\bigl(\bigcup_nA_n\bigr)$
        (continuity from below).
  \item If $A_1\supseteq A_2\supseteq\cdots$, then $\Prob(A_n)\to\Prob\bigl(\bigcap_nA_n\bigr)$
        (continuity from above).
\end{enumerate}
\end{theorem}

\begin{proof}
(1) The sets $B_1 = A_1$, $B_k = A_k\setminus A_{k-1}$ are disjoint, with
$\bigcup_{k\leq n}B_k = A_n$ and $\bigcup_kB_k = \bigcup_nA_n$. By
$\sigma$-additivity, $\Prob(\bigcup_nA_n) = \sum_k\Prob(B_k) = \lim_n\sum_{k\leq n}\Prob(B_k) = \lim_n\Prob(A_n)$.
(2) Apply (1) to the increasing complements $A_n^c$.
\end{proof}

These properties hold for \emph{every} probability measure --- they are
equivalent to $\sigma$-additivity. A monotone sequence always converges, to its
union or intersection, so the theorem is a special case of the next one.

\begin{theorem}[Convergent events]\label{convergence-of-events:thm:convergent}
If $A_n\to A$ in the set-theoretic sense, then $\Prob(A_n)\to\Prob(A)$.
\end{theorem}

\begin{proof}
Convergence means $\bbone_{A_n}\to\bbone_A$ pointwise. The indicators are bounded
by $1$, so the dominated convergence theorem gives
$\Prob(A_n) = \E[\bbone_{A_n}]\to\E[\bbone_A] = \Prob(A)$.
\end{proof}

No ``continuity'' assumption on $\Prob$ (absence of atoms, say) is needed.
Convergence of the events is however more than required: sets of probability
zero should not matter.

\begin{definition}[Convergence in symmetric difference]
$A_n\to A$ \emph{in symmetric difference} if $\Prob(A_n\,\Delta\,A)\to0$, where
$A_n\,\Delta\,A = (A_n\setminus A)\cup(A\setminus A_n)$.
\end{definition}

\begin{proposition}\label{convergence-of-events:prop:symdiff}
If $\Prob(A_n\,\Delta\,A)\to0$, then $\Prob(A_n)\to\Prob(A)$.
\end{proposition}

\begin{proof}
Splitting $A_n$ and $A$ along $A_n\cap A$,
$\Prob(A_n) = \Prob(A_n\cap A) + \Prob(A_n\setminus A)$ and
$\Prob(A) = \Prob(A_n\cap A) + \Prob(A\setminus A_n)$. Subtracting,
$\abs{\Prob(A_n) - \Prob(A)}\leq\Prob(A_n\setminus A) + \Prob(A\setminus A_n) = \Prob(A_n\,\Delta\,A)$.
\end{proof}

Set-theoretic convergence implies convergence in symmetric difference (apply
dominated convergence to $\bbone_{A_n\Delta A} = \abs{\bbone_{A_n} - \bbone_A}$),
but not conversely. Finally, without any convergence, Fatou's lemma applied to
indicators brackets the probabilities.

\begin{proposition}[Fatou bounds]\label{convergence-of-events:prop:fatou}
For any sequence of events,
\[
  \Prob\bigl(\liminf_nA_n\bigr)\leq\liminf_n\Prob(A_n)\leq\limsup_n\Prob(A_n)\leq\Prob\bigl(\limsup_nA_n\bigr).
\]
In particular $\Prob(A_n)$ converges whenever $\Prob(\limsup A_n\setminus\liminf A_n) = 0$.
\end{proposition}

\begin{proof}
The events $C_N = \bigcap_{n\geq N}A_n$ increase to $\liminf A_n$ and satisfy
$C_N\subseteq A_n$ for $n\geq N$, so
$\Prob(\liminf A_n) = \lim_N\Prob(C_N)\leq\liminf_n\Prob(A_n)$. The last inequality
follows by applying this to the complements.
\end{proof}

\section{Examples}\label{convergence-of-events:sec:examples}

Let $\Omega = [0,1]$ with the uniform probability.

\begin{example}[A monotone family]
$A_\epsilon = [0,\frac12 + \epsilon)$ decreases as $\epsilon\downarrow0$, with
intersection $[0,\frac12]$. By continuity from above,
$\Prob(A_\epsilon) = \frac12 + \epsilon\to\frac12$.
\end{example}

\begin{example}[Events that do not converge, probabilities that do]\label{convergence-of-events:ex:oscillate}
Let $A_n = [0,\frac12 + \frac{(-1)^n}{n})$: slightly larger than $[0,\frac12)$
for even $n$, slightly smaller for odd $n$. Every $\omega<\frac12$ is eventually
always in $A_n$ and every $\omega>\frac12$ eventually always out, but
$\omega = \frac12$ belongs to all even-indexed sets and to no odd-indexed one. So
$\liminf A_n = [0,\frac12)$ and $\limsup A_n = [0,\frac12]$: the events do not
converge. Yet $\Prob(A_n) = \frac12 + \frac{(-1)^n}{n}\to\frac12$, because
$\liminf$ and $\limsup$ differ by the null set $\{\frac12\}$
(\cref{convergence-of-events:prop:fatou}); equivalently,
$\Prob(A_n\,\Delta\,[0,\frac12])\leq\frac1n\to0$.
\end{example}

\begin{example}[Probabilities that do not converge]
Let $A_n = [0,\frac12]$ for even $n$ and $A_n = [0,1]$ for odd $n$. Then
$\liminf A_n = [0,\frac12]$, $\limsup A_n = [0,1]$, and $\Prob(A_n)$ alternates
between $\frac12$ and $1$. The events disagree on a set of positive
probability, and nothing can be said.
\end{example}

\section{Application: the symmetrization lemma}\label{convergence-of-events:sec:symmetrization}

In the proof of \cref{symmetrization-lemma:lem:main}, with $Pg$ the true risk,
$P_ng$ and $P'_ng$ the empirical risks on the sample and on the ghost sample,
we met the events
\[
  A = \Bigl\{\sup_{g\in\mathcal G}(Pg - P_ng)\geq t\Bigr\},
  \qquad
  A_\epsilon = \bigl\{Pg_\epsilon - P_ng_\epsilon\geq t-\epsilon\bigr\},
\]
where $g_\epsilon$ is a data-dependent function within $\epsilon$ of the
supremum.

\begin{pitfall}[Do not try to prove $\Prob(A_\epsilon)\to\Prob(A)$]
Because $g_\epsilon$ changes with $\epsilon$, the events $A_\epsilon$ need not be
monotone and need not converge to anything, so none of the theorems above
applies. The proof never needs such a limit. The inclusion $A\subseteq A_\epsilon$
(if the supremum is at least $t$, then $g_\epsilon$ achieves at least $t-\epsilon$)
gives $\Prob(A)\leq\Prob(A_\epsilon)$ for every $\epsilon$ directly. The only limit
concerns the right-hand side,
\[
  \Prob\Bigl[\sup_g(P'_ng - P_ng)\geq\tfrac{t-\epsilon}2\Bigr]
  \xrightarrow[\epsilon\downarrow0]{}
  \Prob\Bigl[\sup_g(P'_ng - P_ng)\geq\tfrac t2\Bigr],
\]
and there the events \emph{decrease} as $\epsilon\downarrow0$ with intersection
$\{\sup_g(P'_ng - P_ng)\geq\frac t2\}$, so continuity from above applies.
\end{pitfall}

\begin{keyidea}
When a proof seems to require a limit of events, look first for an inclusion
($A\subseteq A_\epsilon$) or a monotone family: either gives the inequality
directly, with no convergence theorem to check.
\end{keyidea}

\begin{chaptersummary}
\begin{itemize}[nosep]
  \item $\liminf A_n$ = ``eventually always'', $\limsup A_n$ = ``infinitely
        often''; $A_n\to A$ when both equal $A$, i.e.\ when indicators converge
        pointwise.
  \item For every probability measure: monotone sequences and convergent
        sequences satisfy $\Prob(A_n)\to\Prob(A)$; so does convergence in
        symmetric difference.
  \item Fatou: $\Prob(\liminf A_n)\leq\liminf\Prob(A_n)\leq\limsup\Prob(A_n)\leq\Prob(\limsup A_n)$.
  \item In the symmetrization lemma an inclusion replaces a limit of events, and
        the one limit taken is along a monotone family.
\end{itemize}
\end{chaptersummary}

\section*{Exercises}
\addcontentsline{toc}{section}{Exercises}

\begin{exercise}\label{convergence-of-events:exr:1}
Let $A_n = [0, 1/n]$ on $\Omega = [0,1]$ with the uniform probability. Compute
$\liminf A_n$, $\limsup A_n$ and $\lim\Prob(A_n)$.
\solutionlink{convergence-of-events:exr:1}
\end{exercise}

\begin{exercise}\label{convergence-of-events:exr:2}
(Borel--Cantelli) Show that if $\sum_n\Prob(A_n)<\infty$, then
$\Prob(\limsup_nA_n) = 0$.
\hint{$\limsup A_n\subseteq\bigcup_{n\geq N}A_n$ for every $N$.}
\solutionlink{convergence-of-events:exr:2}
\end{exercise}

\begin{exercise}\label{convergence-of-events:exr:3}
Show that $\bbone_{\limsup A_n} = \limsup_n\bbone_{A_n}$ and
$\bbone_{\liminf A_n} = \liminf_n\bbone_{A_n}$ pointwise.
\solutionlink{convergence-of-events:exr:3}
\end{exercise}
\end{notechapter}
